%-----------------------------------------------
%  Engineer's & Master's Thesis Template
%  Copyleft by Artur M. Brodzki & Piotr Woźniak
%  Warsaw University of Technology, 2019-2022
%-----------------------------------------------

\documentclass[
    bindingoffset=5mm,  % Binding offset
    footnoteindent=3mm, % Footnote indent
    hyphenation=true    % Hyphenation turn on/off
]{src/wut-thesis}

\graphicspath{{tex/img/}} % Katalog z obrazkami.
\addbibresource{bibliografia.bib} % Plik .bib z bibliografią

%-------------------------------------------------------------
% Wybór wydziału:
%  \facultyeiti: Wydział Elektroniki i Technik Informacyjnych
%  \facultymeil: Wydział Mechaniczny Energetyki i Lotnictwa
% --
% Rodzaj pracy: \EngineerThesis, \MasterThesis
% --
% Wybór języka: \langpol, \langeng
%-------------------------------------------------------------
\facultyeiti    % Wydział Elektroniki i Technik Informacyjnych
\MasterThesis % Praca inżynierska
\langpol % Praca w języku polskim

\begin{document}

%------------------
% Strona tytułowa
%------------------
\instytut{Automatyki i Informatyki Stosowanej}
\kierunek{Informatyka}
\specjalnosc{Inteligetne Systemy}
\title{
    Porównanie wydajności i możliwości współczesnych silników gier komputerowych
}
% Title in English for English theses
% In English theses, you may remove this command
\engtitle{
    Comparison of performance and capabilities of modern computer games engines
}

\author{inż. Krzysztof Rudnicki}
\album{307585}
\promotor{dr inż. Michał Chwesiuk}
\date{\the\year}
\maketitle

%-------------------------------------
% Streszczenie po polsku dla \langpol
% English abstract if \langeng is set
%-------------------------------------
\cleardoublepage % Zaczynamy od nieparzystej strony
\abstract \lipsum[1-3]
\keywords XXX, XXX, XXX

%----------------------------------------
% Streszczenie po angielsku dla \langpol
% Polish abstract if \langeng is set
%----------------------------------------
\clearpage
\secondabstract \kant[1-3]
\secondkeywords XXX, XXX, XXX

\pagestyle{plain}

%--------------
% Spis treści
%--------------
\cleardoublepage % Zaczynamy od nieparzystej strony
\tableofcontents

%------------
% Rozdziały
%------------
\cleardoublepage % Zaczynamy od nieparzystej strony
\pagestyle{headings}

\clearpage % Rozdziały zaczynamy od nowej strony.
\section{Praefatio}

% Akapit z cytatem
\lipsum[1] \cite{goossens93}

% Przykładowy obrazek
\begin{figure}[!h]
    % Wyrównanie obrazka, szerokość i plik
    % Zamiast width można też użyć height, etc.
    \centering \includegraphics[width=0.5\linewidth]{logopw.png}
    % Podpis umieszczamy pod obrazkiem
    % znacznik \caption służy również do wygenerowania numeru obrazka
    \caption{Tradycyjne godło Politechniki Warszawskiej}
    % \label pozwala odwołać się do obrazka w innych miejscach za pomocą \ref
    % odwołanie \ref renderuje się jako numer obrazka,
    % dlatego zawsze najpierw używaj \caption a potem \label
    \label{fig:tradycyjne-logo-pw}
\end{figure}

% \ref wyrenderuje się jako 'Reference to image 1.1'
\lipsum[2] Reference to image \ref{fig:tradycyjne-logo-pw}.

% Lista punktowana
% Parametr label ustawia symbol punktora
\begin{itemize}
    \item Item 1:
    \begin{itemize}[label=---]
        \item item 1.1;
        \item item 1.2;
    \end{itemize}
    \item Item 2;
    \item Item 3.
\end{itemize}

\lipsum[3]

% Lista numerowana w formacie 1.a).ii
% Tutaj również można stosować \label
\begin{enumerate}
    \item Item 1:
    \begin{enumerate}
        \item item 1.1;
        \item item 1.2:
        \begin{enumerate}
            \item item 1.2.1;
            \item item 1.2.2;
        \end{enumerate}
        \item item 1.3;
    \end{enumerate}
    \item Item 2;
    \item Item 3.
\end{enumerate}

% Przypis dolny \footnote
\lipsum[4] Lorem ipsum dolor sit amet\footnote{Lorem ipsum dolor sit amet, consectetur adipiscing elit, sed do eiusmod tempor incididunt ut labore et dolore magna aliqua. Ut enim ad minim veniam, quis nostrud exercitation ullamco laboris nisi ut aliquip ex ea commodo consequat.}, consectetur adipiscing elit.

% Przykładowa tabela: wyśrodkowana i renderowana
% w miejscu wstawienia: !h = !h[ere]
% Domyślnie tabele trafiają na górę strony
\begin{table}[!h] \centering
    % Podpis tabeli umieszczamy od góry
    \caption{Przykładowa tabela.}
    \label{tab:tabela1}

    % Tabela z trzema kolumnami:
    % dwie wyrównanie do środka [c], a ostatnia do prawej [r]
    % szerokość kolumn automatyczna (równa szerokości tekstu)
    \begin{tabular}{| c | c | r |} \hline
        Kolumna 1       & Kolumna 2 & Liczba \\ \hline\hline
        cell1           & cell2     & 60     \\ \hline
        cell4           & cell5     & 43     \\ \hline
        cell7           & cell8     & 20,45  \\ \hline
        % Komórka o szerokości dwóch kolumn, wyrównana do prawej
        % Przypisy dolne w tabelach wstawiamy przez \tablefootnote
        \multicolumn{2}{|r|}{Suma\tablefootnote{Table footnote.}} & 123,45 \\ \hline
    \end{tabular}

\end{table}

Lorem ipsum dolor sit amet.
                    % Wstęp
\clearpage
\section{Przegląd literatury i istniejących rozwiązań}

\subsection{Historia rozwoju silników gier}
W tym podrozdziale zostanie przedstawiona ewolucja silników gier od prostych bibliotek graficznych po współczesne, kompleksowe środowiska deweloperskie.

\subsection{Klasyfikacja silników gier}
\subsubsection{Silniki komercyjne vs. open source}
% Analiza różnic między rozwiązaniami komercyjnymi a otwartymi

\subsubsection{Silniki specjalistyczne vs. uniwersalne}
% Porównanie silników dedykowanych konkretnym gatunkom gier z rozwiązaniami uniwersalnymi

\subsection{Aktualny stan badań}
Przegląd istniejących prac naukowych dotyczących porównania silników gier, ze szczególnym uwzględnieniem metodologii badawczych i uzyskanych wyników.

\subsection{Identyfikacja luk badawczych}
Wskazanie obszarów, które nie zostały dotychczas wystarczająco zbadane oraz uzasadnienie potrzeby przeprowadzenia niniejszych badań.

% Bibliografia dla tego rozdziału będzie zawierać odniesienia do:
% - publikacji naukowych
% - dokumentacji technicznej
% - raportów branżowych
% - analiz rynkowych
       % Przegląd literatury i istniejących rozwiązań
\clearpage
\section{Charakterystyka współczesnych silników gier}

\subsection{Kryteria wyboru silników do analizy}
W ramach niniejszej pracy wybrano reprezentatywne silniki gier reprezentujące różne segmenty rynku:
\begin{itemize}
    \item Popularność i rozpowszechnienie
    \item Różnorodność technologiczna
    \item Wsparcie dla różnych platform
    \item Dostępność dokumentacji i narzędzi
\end{itemize}

\subsection{Unity}
\subsubsection{Architektura i technologie}
% Opis architektury silnika Unity, stosowanych technologii

\subsubsection{Możliwości i funkcjonalności}
% Szczegółowy opis oferowanych funkcji

\subsubsection{Narzędzia deweloperskie}
% Edytor, debugger, profiler, etc.

\subsection{Unreal Engine}
\subsubsection{Architektura i technologie}
% Opis architektury Unreal Engine

\subsubsection{Możliwości i funkcjonalności}
% Blueprint system, C++ support, rendering pipeline

\subsubsection{Narzędzia deweloperskie}
% Unreal Editor, Visual Scripting, etc.

\subsection{Godot}
\subsubsection{Architektura i technologie}
% Opis architektury silnika Godot

\subsubsection{Możliwości i funkcjonalności}
% GDScript, scene system, node-based architecture

\subsubsection{Narzędzia deweloperskie}
% Integrated editor, scripting tools

\subsection{Inne analizowane silniki}
% Krótka charakterystyka dodatkowych silników (np. CryEngine, GameMaker Studio)

\subsection{Porównanie tabelaryczne podstawowych cech}
% Tabela porównawcza kluczowych parametrów
              % Charakterystyka współczesnych silników gier
\clearpage
\section{Metodologia badań i kryteria porównania}

\subsection{Założenia metodologiczne}
\subsubsection{Cel badań}
Głównym celem badań jest obiektywne porównanie wydajności i możliwości wybranych silników gier w kontrolowanych warunkach.

\subsubsection{Hipotezy badawcze}
\begin{enumerate}
    \item Silniki komercyjne oferują lepszą wydajność niż rozwiązania open source
    \item Kompleksowość funkcjonalności wpływa negatywnie na wydajność
    \item Łatwość użycia jest odwrotnie proporcjonalna do możliwości konfiguracji
\end{enumerate}

\subsection{Kryteria porównania}
\subsubsection{Wydajność}
\begin{itemize}
    \item Szybkość renderowania (FPS)
    \item Zużycie pamięci RAM
    \item Obciążenie procesora
    \item Zużycie pamięci karty graficznej
    \item Czas ładowania scen
\end{itemize}

\subsubsection{Funkcjonalność}
\begin{itemize}
    \item Wsparcie dla różnych typów renderingu
    \item Systemy fizyki
    \item Systemy audio
    \item Wsparcie dla VR/AR
    \item Możliwości skryptowania
\end{itemize}

\subsubsection{Użyteczność}
\begin{itemize}
    \item Intuicyjność interfejsu
    \item Jakość dokumentacji
    \item Dostępność tutoriali
    \item Wsparcie społeczności
    \item Czas potrzebny na naukę
\end{itemize}

\subsection{Środowisko testowe}
\subsubsection{Specyfikacja sprzętowa}
% Szczegółowy opis konfiguracji sprzętowej używanej do testów

\subsubsection{Specyfikacja oprogramowania}
% Wersje systemów operacyjnych, sterowników, silników

\subsection{Projekt testów}
\subsubsection{Scenariusze testowe}
% Opis konkretnych scenariuszy używanych do porównania

\subsubsection{Metryki i wskaźniki}
% Definicja mierzonych parametrów i sposobów ich pomiaru
               % Metodologia badań i kryteria porównania
\clearpage
\section{Testy wydajności}

\subsection{Metodyka przeprowadzania testów}
\subsubsection{Przygotowanie środowiska testowego}
% Opis procesu przygotowania spójnego środowiska dla wszystkich silników

\subsubsection{Standaryzacja warunków testowych}
% Zapewnienie porównywalności wyników między różnymi silnikami

\subsection{Test renderowania 2D}
\subsubsection{Założenia testu}
% Opis scenariusza testowego dla grafiki 2D

\subsubsection{Wyniki pomiarów}
% Tabele i wykresy przedstawiające wyniki dla każdego silnika

\subsubsection{Analiza wyników}
% Interpretacja uzyskanych danych

\subsection{Test renderowania 3D}
\subsubsection{Scenariusz podstawowy}
% Test renderowania prostej sceny 3D

\subsubsection{Scenariusz zaawansowany}
% Test złożonych scen z wieloma obiektami, oświetleniem, cieniami

\subsubsection{Porównanie wyników}
% Analiza porównawcza wydajności renderowania 3D

\subsection{Test systemów fizyki}
\subsubsection{Symulacja kolizji}
% Testy wydajności przy różnej liczbie obiektów fizycznych

\subsubsection{Symulacja płynów i cząstek}
% Testy zaawansowanych systemów fizyki

\subsection{Test zużycia zasobów systemowych}
\subsubsection{Zużycie pamięci RAM}
% Pomiary zużycia pamięci operacyjnej

\subsubsection{Obciążenie procesora}
% Analiza wykorzystania CPU

\subsubsection{Zużycie pamięci GPU}
% Pomiary zużycia pamięci karty graficznej

\subsection{Test wydajności na różnych platformach}
\subsubsection{Testy na PC (Windows/Linux)}
% Porównanie wydajności na systemach desktop

\subsubsection{Testy na urządzeniach mobilnych}
% Analiza wydajności na platformach mobilnych (jeśli dotyczy)

\subsection{Podsumowanie wyników testów wydajności}
% Zestawienie wszystkich uzyskanych wyników w formie tabelarycznej i graficznej
          % Testy wydajności
\clearpage
\section{Analiza możliwości i funkcjonalności}

\subsection{Metodyka oceny funkcjonalności}
\subsubsection{Kryteria oceny}
% Opis systemu punktowego lub skali ocen

\subsubsection{Proces walidacji}
% Sposób weryfikacji dostępności i jakości funkcji

\subsection{Analiza możliwości renderingu}
\subsubsection{Wsparcie dla różnych technik renderingu}
% Forward rendering, deferred rendering, ray tracing, etc.

\subsubsection{Systemy materiałów i shaderów}
% Możliwości tworzenia i edycji materiałów

\subsubsection{Systemy oświetlenia}
% Real-time lighting, baked lighting, global illumination

\subsection{Systemy fizyki i symulacji}
\subsubsection{Rigid body physics}
% Analiza możliwości symulacji ciał sztywnych

\subsubsection{Soft body physics}
% Wsparcie dla symulacji ciał miękkich

\subsubsection{Systemy cząstek}
% Możliwości tworzenia efektów cząsteczkowych

\subsection{Systemy audio}
\subsubsection{Wsparcie formatów audio}
% Obsługiwane formaty i kodeki

\subsubsection{Przestrzenny dźwięk 3D}
% Możliwości pozycjonowania dźwięku w przestrzeni

\subsubsection{Efekty audio i DSP}
% Dostępne efekty i systemy przetwarzania dźwięku

\subsection{Narzędzia deweloperskie}
\subsubsection{Edytory wizualne}
% Analiza jakości i funkcjonalności edytorów

\subsubsection{Systemy debugowania}
% Dostępne narzędzia do debugowania

\subsubsection{Profilowanie wydajności}
% Wbudowane narzędzia do analizy wydajności

\subsection{Wsparcie dla platform docelowych}
\subsubsection{Platformy desktop}
% Windows, macOS, Linux

\subsubsection{Platformy mobilne}
% iOS, Android

\subsubsection{Konsole}
% PlayStation, Xbox, Nintendo Switch

\subsubsection{Platformy VR/AR}
% Wsparcie dla rzeczywistości wirtualnej i rozszerzonej

\subsection{Ekosystem i rozszerzalność}
\subsubsection{Asset Store / Marketplace}
% Dostępność gotowych zasobów i rozszerzeń

\subsubsection{Wsparcie społeczności}
% Aktywność społeczności, dostępność pomocy

\subsubsection{Dokumentacja i materiały edukacyjne}
% Jakość oficjalnej dokumentacji i tutoriali
        % Analiza możliwości i funkcjonalności
\clearpage
\section{Porównanie wyników i analiza}

\subsection{Synteza wyników badań}
\subsubsection{Zestawienie wyników testów wydajności}
% Kompleksowe podsumowanie wszystkich pomiarów wydajności

\subsubsection{Zestawienie analizy funkcjonalności}
% Podsumowanie oceny możliwości każdego silnika

\subsection{Analiza wielokryterialna}
\subsubsection{Macierz porównawcza}
% Tabela porównująca wszystkie silniki według wszystkich kryteriów

\subsubsection{Analiza wag kryteriów}
% Określenie ważności poszczególnych aspektów dla różnych typów projektów

\subsection{Przypadki użycia}
\subsubsection{Gry indie}
% Rekomendacje dla małych zespołów deweloperskich

\subsubsection{Gry mobilne}
% Najlepsze rozwiązania dla platform mobilnych

\subsubsection{Gry AAA}
% Analiza przydatności dla dużych projektów komercyjnych

\subsubsection{Gry VR/AR}
% Specyficzne wymagania dla rzeczywistości wirtualnej

\subsection{Analiza korelacji}
\subsubsection{Związek między wydajnością a funkcjonalnością}
% Analiza statystyczna zależności między różnymi parametrami

\subsubsection{Wpływ złożoności na użyteczność}
% Badanie relacji między liczbą funkcji a łatwością użycia

\subsection{Ograniczenia badań}
\subsubsection{Ograniczenia metodologiczne}
% Omówienie ograniczeń zastosowanej metodologii

\subsubsection{Ograniczenia techniczne}
% Wpływ środowiska testowego na wyniki

\subsubsection{Ograniczenia czasowe}
% Wpływ dynamicznego rozwoju silników na aktualność wyników

\subsection{Weryfikacja hipotez badawczych}
% Odniesienie uzyskanych wyników do postawionych na początku hipotez

\subsection{Implikacje praktyczne}
% Wnioski dla praktyki tworzenia gier i wyboru silnika
        % Porównanie wyników i analiza
\clearpage
\section{Podsumowanie i wnioski}

\subsection{Główne wyniki badań}
\subsubsection{Odpowiedzi na pytania badawcze}
% Zwięzłe odpowiedzi na postawione na początku pytania

\subsubsection{Weryfikacja hipotez}
% Potwierdzenie lub odrzucenie hipotez badawczych

\subsection{Wnioski praktyczne}
\subsubsection{Rekomendacje dla deweloperów}
% Praktyczne wskazówki dotyczące wyboru silnika gier

\subsubsection{Wytyczne dla różnych typów projektów}
% Szczegółowe rekomendacje w zależności od rodzaju gry

\subsection{Wkład naukowy}
\subsubsection{Nowatorskie aspekty badań}
% Elementy które wnoszą nową wiedzę do dziedziny

\subsubsection{Znaczenie dla branży}
% Potencjalny wpływ wyników na przemysł gier

\subsection{Ograniczenia i przyszłe badania}
\subsubsection{Identyfikacja ograniczeń}
% Szczere omówienie ograniczeń przeprowadzonych badań

\subsubsection{Propozycje dalszych badań}
% Sugestie dla przyszłych prac badawczych w tej dziedzinie

\subsubsection{Rozwój metodologii}
% Możliwości udoskonalenia zastosowanej metodologii

\subsection{Refleksje końcowe}
% Osobiste refleksje autora na temat przeprowadzonych badań i ich znaczenia

\subsection{Znaczenie wyników w kontekście rozwoju technologii}
% Umiejscowienie wyników w szerszym kontekście rozwoju technologii gier

% Ten rozdział powinien stanowić logiczne zamknięcie całej pracy
% i pokazać, że cele zostały osiągnięte
              % Podsumowanie i wnioski

%---------------
% Bibliografia
%---------------
\cleardoublepage % Zaczynamy od nieparzystej strony
\printbibliography
\clearpage

% Wykaz symboli i skrótów.
% Pamiętaj, żeby posortować symbole alfabetycznie
% we własnym zakresie. Makro \acronymlist
% generuje właściwy tytuł sekcji, w zależności od języka.
% Makro \acronym dodaje skrót/symbol do listy,
% zapewniając podstawowe formatowanie.
\acronymlist
\acronym{EiTI}{Wydział Elektroniki i Technik Informacyjnych}
\acronym{PW}{Politechnika Warszawska}
\acronym{WEIRD}{ang. \emph{Western, Educated, Industrialized, Rich and Democratic}}
\vspace{0.8cm}

%--------------------------------------
% Spisy: rysunków, tabel, załączników
%--------------------------------------
\pagestyle{plain}

\listoffigurestoc    % Spis rysunków.
\vspace{1cm}         % vertical space
\listoftablestoc     % Spis tabel.
\vspace{1cm}         % vertical space
\listofappendicestoc % Spis załączników

%-------------
% Załączniki
%-------------

% Obrazki i tabele w załącznikach nie trafiają do spisów
\captionsetup[figure]{list=no}
\captionsetup[table]{list=no}

% Załącznik 1
\clearpage
\appendix{Nazwa załącznika 1}
\lipsum[1-3]
\begin{figure}[!h]
	\centering \includegraphics[width=0.5\linewidth]{logopw2.png}
	\caption{Obrazek w załączniku.}
\end{figure}
\lipsum[4-7]

% Załącznik 2
\clearpage
\appendix{Nazwa załącznika 2}
\lipsum[1-2]
\begin{table}[!h] \centering
    \caption{Tabela w załączniku.}
    \begin{tabular} {| c | c | r |} \hline
        Kolumna 1       & Kolumna 2 & Liczba \\ \hline\hline
        cell1           & cell2     & 60     \\ \hline
        \multicolumn{2}{|r|}{Suma:} & 123,45 \\ \hline
    \end{tabular}
\end{table}
\lipsum[3-4]

% Używając powyższych spisów jako szablonu,
% możesz dodać również swój własny wykaz,
% np. spis algorytmów.

\end{document} % Dobranoc.
